\section{Main Theorem Assembly}
\label{sec:main-assembly}

We now assemble the declarations required by \code{Jacobian/Challenge.lean}.

\begin{theorem}[Challenge API from the analytic construction]
\label{thm:challenge-api}
\uses{thm:fd-holomorphic-one-forms,thm:genus-zero-classification,def:analytic-jacobian,def:abel-jacobi,lem:aj-self,lem:aj-holomorphic,thm:abel-jacobi-injective,thm:push-pull-functoriality,thm:pushforward-pullback}
The analytic period-lattice construction supplies all definitions and theorems
in the Jacobian Challenge API.
\lean{genus, genus_eq_zero_iff_homeo, Jacobian, Jacobian.ofCurve, Jacobian.pushforward, Jacobian.pullback, Jacobian.pushforward_pullback}
\leanok
\end{theorem}

\begin{layman}
This is the moment all the bookkeeping pays off. We assemble the public
API of the Jacobian Challenge — the small set of definitions and
theorems that \texttt{Jacobian/Challenge.lean} exposes — directly from
the analytic period-lattice construction. \emph{Genus} is the
finite-dimensional dimension of holomorphic 1-forms; the genus-zero
classification matches it to the topological sphere. The
\emph{Jacobian} is the dual-of-1-forms quotient by the period lattice,
inheriting all its algebraic, topological, and analytic structure from
the complex-torus package. \emph{Abel–Jacobi} is the path-integration
map, holomorphic and (in positive genus) injective. \emph{Pullback}
and \emph{pushforward} between Jacobians come from the dual of the
1-form maps, and the trace identity gives the key
``push-then-pull equals multiplication-by-degree'' equation. Why
bother? This is the deliverable. Once every constituent theorem is
formally proved in Lean, this assembly node certifies that the entire
\texttt{Challenge.lean} API holds — finishing the challenge.
\end{layman}

\begin{proof}
Define \code{genus X} to be the analytic genus
\(\dim_{\C}H^{0}(X,\Omega^{1}_{X})\). Theorem~\ref{thm:fd-holomorphic-one-forms}
ensures this dimension is finite, and Theorem~\ref{thm:genus-zero-classification}
proves \code{genus_eq_zero_iff_homeo}.

Define \code{Jacobian X} to be the complex torus
\(\HdualX/\periods_X\). Theorem~\ref{thm:period-lattice} makes
\(\periods_X\) a full lattice, and Proposition~\ref{prop:complex-torus-package}
gives the additive group, topology, Hausdorffness, compactness, charted-space,
manifold, and Lie additive group structures.

Define \code{Jacobian.ofCurve P} to be \(\AJ_P\). Lemma~\ref{lem:aj-self}
proves \code{ofCurve_self}; Lemma~\ref{lem:aj-holomorphic} proves
\code{ofCurve_contMDiff}; Theorem~\ref{thm:abel-jacobi-injective} proves
\code{ofCurve_inj}.

For a holomorphic map \(f:X\to Y\), define \code{ContMDiff.degree f hf} as
the analytic degree. Define \code{Jacobian.pullback f hf} and
\code{Jacobian.pushforward f hf} by the descended maps from
Lemma~\ref{lem:push-pull-descend}. The holomorphicity, identity, and
composition laws are Theorem~\ref{thm:push-pull-functoriality}.

Finally, \code{pushforward_pullback} is exactly
Theorem~\ref{thm:pushforward-pullback}.
\end{proof}
